\chapter*{Заключение}                       % Заголовок
\addcontentsline{toc}{chapter}{Заключение}  % Добавляем его в оглавление

%% Согласно ГОСТ Р 7.0.11-2011:
%% 5.3.3 В заключении диссертации излагают итоги выполненного исследования, рекомендации, перспективы дальнейшей разработки темы.
%% 9.2.3 В заключении автореферата диссертации излагают итоги данного исследования, рекомендации и перспективы дальнейшей разработки темы.
%% Поэтому имеет смысл сделать эту часть общей и загрузить из одного файла в автореферат и в диссертацию:

Основные результаты работы заключаются в следующем.
%% Согласно ГОСТ Р 7.0.11-2011:
%% 5.3.3 В заключении диссертации излагают итоги выполненного исследования, рекомендации, перспективы дальнейшей разработки темы.
%% 9.2.3 В заключении автореферата диссертации излагают итоги данного исследования, рекомендации и перспективы дальнейшей разработки темы.
\begin{enumerate}
  \item На основе анализа \ldots
  \item Численные исследования показали, что \ldots
  \item Математическое моделирование показало \ldots
  \item Для выполнения поставленных задач был создан \ldots
\end{enumerate}


Основные результаты работы заключаются в следующем.
1. Произведен обзор методов переноса знаний при решении проблем обра
ботки естественного языка.
2. Произведено внутреннее и внешнее сравнение контекстно-независи
мых векторных представлений слов, обученных разными способами
на ограниченной обучающей выборке для трех малоресурсных языков.
Обученные векторные представления выложены в открытый доступ.
3. Разработана вопросно-ответная модель на основе BERT с применением
библиотеки transformers от HuggingFace.
4. Показано, что межъязыковой перенос позволяет использовать обще
доступные обучающие данные английского языка при дообучении
M-BERT для вопросно-ответной задачи, таким образом, отсутствует
необходимость в сборе полноценного обучающего датасета целевого
языка.
5. Экспериментально установлено, что применяя метод межъязыково
го переноса, M-BERT, дообученный с применением многоязычной
обучающей выборки, имеет сопоставимое качество с языко-специ
фичными BERT для вопросно-ответной задачи, тем самым отпадает
необходимость в вычислительных ресурсах для предобучения языко
специфичных BERT.
6. Разработана оригинальная модель GOLOMB, которая использует во
просно-ответную модель SQuAD для отслеживания состояния диалога.
7. Модели, обученные в рамках работы, доступны в открытом доступе.
% И какая-нибудь заключающая фраза.
В данной диссертации проведено комплексное исследование нейронных сетей с памятью, охватывающее как обзор ключевых архитектур, так и разработку новых подходов к обучению и обработке длинных контекстов. Представлены концептуальные основы для анализа линейных и нелинейных моделей обработки последовательностей, предложена архитектура Recurrent Memory Transformer (RMT), которая сочетает в себе рекуррентность и память для повышения эффективности обработки данных. Разработанный бенчмарк BABILong позволил выявить слабые стороны современных языковых моделей при работе с длинными контекстами и продемонстрировал потенциал рекуррентных трансформеров в этой области. Эти результаты закладывают основу для дальнейших исследований интеграции памяти в нейронные сети и их применения к задачам, требующим долгосрочных зависимостей и сложного рассуждения.

% Последний параграф может включать благодарности.  В заключение автор
% выражает благодарность и большую признательность научному руководителю
% Иванову~И.\,И. за поддержку, помощь, обсуждение результатов и~научное
% руководство. Также автор благодарит Сидорова~А.\,А. и~Петрова~Б.\,Б.
% за помощь в~работе с~образцами, Рабиновича~В.\,В. за предоставленные
% образцы и~обсуждение результатов, Занудятину~Г.\,Г. и авторов шаблона
% *Russian-Phd-LaTeX-Dissertation-Template* за~помощь в оформлении
% диссертации. Автор также благодарит много разных людей
% и~всех, кто сделал настоящую работу автора возможной.
